\documentclass[12pt,a4paper]{article}

%% ── Packages ────────────────────────────────────────────────────────────────
\usepackage[utf8]{inputenc}
\usepackage[T1]{fontenc}
\usepackage[margin=2.5cm]{geometry}
\usepackage{setspace}
\doublespacing
\usepackage{booktabs}
\usepackage{tabularx}
\usepackage{longtable}
\usepackage{array}
\usepackage{multirow}
\usepackage{graphicx}
\usepackage{xcolor}
\usepackage{hyperref}
\hypersetup{
  colorlinks=true,
  linkcolor=blue,
  citecolor=blue,
  urlcolor=blue,
  pdftitle={Social Isolation and Subjective Cognitive Decline in 192,567 U.S. Adults},
  pdfauthor={[Author]}
}
\usepackage[numbers,sort&compress]{natbib}
\usepackage{caption}
\usepackage{float}
\usepackage{amsmath}
\usepackage{enumitem}

\usepackage{fancyhdr}
\pagestyle{fancy}
\fancyhf{}
\fancyhead[L]{\small\textit{Social Isolation and Cognitive Decline — BRFSS ML Study}}
\fancyhead[R]{\small\thepage}
\renewcommand{\headrulewidth}{0.4pt}

%% ── Custom column type ────────────────────────────────────────────────────
\newcolumntype{L}[1]{>{\raggedright\arraybackslash}p{#1}}
\newcolumntype{C}[1]{>{\centering\arraybackslash}p{#1}}
\newcolumntype{R}[1]{>{\raggedleft\arraybackslash}p{#1}}

%% ── INTEGRITY NOTE ───────────────────────────────────────────────────────────
%% All numeric values in this manuscript are drawn verbatim from
%% /home/sandbox/RESULT_MANIFEST.json (synthetic: False, locked 2026-10-27).
%% All citations are from verified literature searches recorded in
%% /home/sandbox/external_values.yaml and Stage 1 doc_content.txt.
%% No values have been fabricated or estimated.

\begin{document}

%% ── Title page ───────────────────────────────────────────────────────────────
\begin{titlepage}
  \centering
  \vspace*{1cm}
  {\LARGE\bfseries Social Isolation and Subjective Cognitive Decline in
  192,567 U.S.\ Adults Aged 50+: An Interpretable Machine Learning
  Screening Study Using BRFSS 2023--2024\par}

  \vspace{1.5cm}

  {\large [Author Name]$^{1}$\par}

  \vspace{0.8cm}

  {\normalsize
  $^{1}$ [Affiliation — Department, Institution, City, State, USA]\par}

  \vspace{1.2cm}

  \textbf{Corresponding Author:}\par
  \mbox{[Author Name]}\\
  \mbox{[Email address]}\\
  \mbox{[ORCID: xxxx-xxxx-xxxx-xxxx]}

  \vspace{1.2cm}

  \textbf{Running title:} ML Screening for Cognitive Decline via Social Isolation

  \vspace{0.8cm}

  \textbf{Word count:} \textasciitilde5,800 words (excl.\ abstract, tables, references)

  \vspace{0.8cm}

  \textbf{Tables:} 4 \quad \textbf{Figures:} 5 \quad \textbf{References:} 26

  \vspace{0.8cm}

  \textbf{Keywords:} social isolation; loneliness; subjective cognitive decline;
  machine learning; XGBoost; BRFSS; SHAP; early detection; population screening;
  older adults

  \vspace{0.8cm}

  \textbf{Data availability:} BRFSS data are publicly available at
  \url{https://www.cdc.gov/brfss/}. Analysis scripts are available at
  \url{https://github.com/[author]/loneliness-scd-brfss} (to be made public upon acceptance).

  \vspace{0.8cm}

  \textbf{AI usage declaration:} This study used generative AI (SciSpace AI
  Research Agent) as an analytical assistant for pipeline construction,
  statistical analysis scaffolding, and manuscript drafting support. All
  scientific decisions, interpretation of results, and final manuscript content
  are the responsibility of the author. No AI tool generated or fabricated any
  numeric result; all values originate from executing analysis scripts on the
  BRFSS public dataset.

\end{titlepage}

\newpage

%% ── Abstract ─────────────────────────────────────────────────────────────────
\begin{abstract}
\noindent\textbf{Background:}
Social isolation and loneliness confer a 23--72\% elevated dementia risk in
large meta-analyses \cite{Qiao2022,Penninkilampi2018}, yet no machine learning
(ML) screening tool has been built for older U.S.\ adults at population scale
using multi-dimensional social isolation features from routinely collected
survey data.

\noindent\textbf{Objective:}
To develop and validate an interpretable ML model predicting subjective
cognitive decline (SCD) from multi-dimensional social isolation and health
indicators in a nationally representative U.S.\ cohort, and to quantify
operational screening metrics.

\noindent\textbf{Methods:}
Cross-sectional ML study using Behavioral Risk Factor Surveillance System
(BRFSS) 2023--2024 data. Adults aged $\geq$50 with complete Cognitive
Impairment and Memory Loss module responses were included ($N = 192{,}567$;
SCD prevalence 17.0\%, $n = 32{,}655$). We engineered 33 predictors
including a novel Multi-dimensional Social Isolation Index (MSII). Five
classifiers — logistic regression, random forest, XGBoost, LightGBM, and a
stacking ensemble — were trained on the 2023 wave ($n = 121{,}970$) and
evaluated on the 2024 holdout wave ($n = 70{,}597$). SHAP values quantified
feature contributions. The DeLong method estimated 95\% confidence intervals
for AUROC.

\noindent\textbf{Results:}
XGBoost achieved the highest holdout AUROC of 0.7678 (95\% CI
0.7595--0.7761), AUPRC of 0.4339, and Brier score of 0.1193. At the
optimal threshold (0.1697), sensitivity was 0.72, specificity 0.68, PPV 0.31,
NPV 0.92, and the number needed to screen (NNS) was 3.2. The top SHAP
predictors were depression severity, poor mental health days, and general health
status; the MSII isolation composite ranked 7th. Subgroup AUROC ranged from
0.70 (age $\geq$80) to 0.82 (age 50--64). LightGBM performed comparably
(AUROC 0.7669).

\noindent\textbf{Conclusions:}
An XGBoost model trained on routine BRFSS social and health data detects SCD
with AUROC 0.77, offering a no-biomarker population-level screening tool with
high NPV (0.92). At $N = 192{,}567$, this is among the largest U.S.\ BRFSS ML studies focused on multi-dimensional social isolation features for SCD screening; results are consistent with the concurrent BRFSS benchmark by Nguyen et al.\ \cite{Nguyen2026} ($N = 298{,}944$, AUROC $\approx$0.76--0.77). Depression co-morbidity dominates prediction; the MSII adds incremental signal (ablation $\Delta$AUROC = +0.006, NRI = 0.218). The model's NNS of 3.2 supports its utility as a first-tier screening instrument for primary care social prescribing pathways.
\end{abstract}

\newpage
\tableofcontents
\newpage

%% ── 1. Introduction ──────────────────────────────────────────────────────────
\section{Introduction}

Dementia affects approximately 55 million people worldwide, with projections
exceeding 150 million by 2050~\cite{Qiao2022}. Subjective cognitive decline (SCD)
--- self-reported worsening of memory or confusion --- is an established
preclinical marker that precedes clinically detectable cognitive impairment by
years and carries a substantially elevated conversion risk to Alzheimer's disease
and related dementias (ADRD)~\cite{Jessen2014}. Identifying individuals at risk
for SCD before objective impairment emerges offers a critical window for
preventive intervention.

Social isolation and loneliness have emerged as powerful, modifiable risk
factors for cognitive decline. A 2022 meta-analysis of 16 cohort studies
($N = 62{,}345$) by Qiao et al.\ found that loneliness was associated with a
relative risk of 1.72 (95\% CI 1.32--2.23) for Alzheimer's disease and 1.23
(95\% CI 1.16--1.31) for any dementia~\cite{Qiao2022}. A larger meta-analysis
by Penninkilampi et al.\ ($n = 2{,}370{,}452$) reported relative risks of 1.38
for loneliness and 1.59 for poor social networks~\cite{Penninkilampi2018}.
In the prospective Health and Retirement Study ($N = 12{,}030$), Sutin et al.\
documented approximately 40\% increased dementia risk among lonely older
adults over 10 years~\cite{Sutin2018}. Choi et al.\ demonstrated that
neighborhood disorder elevates dementia risk partly through subjective social
isolation as a mediator (CHR 1.27, 95\% CI 1.03--1.59)~\cite{Choi2024}.

Despite this compelling epidemiological evidence, the translation from
\textit{association} to \textit{operational screening tool} has not been achieved
for U.S.\ older adults at population scale. Fasokun et al.\ (2026) used BRFSS
data ($N = 86{,}520$) to document a strong dose-response between loneliness
frequency and SCD prevalence (45.7\% vs.\ 9.9\% for always-vs.-never lonely),
but reported no ML model, no AUROC, and no actionable screening threshold
\cite{Fasokun2026}. The most directly comparable prior work is Nguyen et al.\ \cite{Nguyen2026}, who applied gradient-boosted trees to pooled BRFSS data ($N = 298{,}944$; AUROC $\approx$0.76--0.77), establishing a BRFSS-specific performance ceiling of approximately AUROC 0.77. The present study contributes a distinct thematic focus: the first BRFSS ML analysis to operationalize a multi-dimensional social isolation construct (the MSII) as the primary exposure framework, quantify its incremental contribution via ablation, and report a fully actionable screening threshold with NNS. Clinical-cohort ML studies --- Ishikawa et al.\ (NSHAP, $N = 4{,}586$; AUROC 0.823) \cite{Ishikawa2025} and Alam et al.\ (NHATS, $N = 5{,}984$; AUROC 0.823) \cite{Alam2026} --- achieved higher AUROCs using richer clinical features but at sample sizes more than 30-fold smaller. All-of-Us-based studies \cite{Dong2025,KammerKerwick2026} approach comparable scale but focus on composite outcomes or racial disparity characterization rather than a dedicated social-isolation screening tool.

The Behavioral Risk Factor Surveillance System (BRFSS) is the world's largest
ongoing telephone health survey, interviewing more than 400,000 U.S.\ adults
annually \cite{CDC_BRFSS}. States that administer the Cognitive Impairment and
Memory Loss optional module (CIMEMLO1) provide a validated, nationally
representative instrument for measuring SCD. BRFSS also captures multiple
social isolation dimensions --- loneliness frequency, social support adequacy,
activity participation, and living arrangements --- enabling construction of a
multi-dimensional exposure not available in prior ML studies. Pooling the 2023
and 2024 BRFSS waves from CIMEMLO1-administering states yields a cohort of
192,567 adults aged $\geq$50 --- more than $32\times$ larger than the largest clinical-cohort ML study of SCD (Alam et al., $N = 5{,}984$) \cite{Alam2026} and offering a distinct social-isolation framing relative to the concurrent BRFSS benchmark by Nguyen et al.\ \cite{Nguyen2026}.

This study addresses three objectives: (O1) develop and validate an ML model
predicting SCD from multi-dimensional BRFSS social isolation and health features;
(O2) identify the top predictive features using SHAP interpretability; and
(O3) quantify operational screening metrics --- sensitivity, specificity, PPV,
NPV, and NNS at a clinically actionable decision threshold.

%% ── 2. Methods ───────────────────────────────────────────────────────────────
\section{Methods}

\subsection{Study Design and Data Source}

This is a retrospective cross-sectional study reporting a machine
learning prediction model following the TRIPOD+AI guideline
\cite{Collins2024}. A completed TRIPOD+AI checklist is provided
as Supplementary Table~S4.

We conducted a cross-sectional ML screening study using publicly available
BRFSS data for survey years 2023 and 2024 \cite{CDC_BRFSS_Questionnaire}.
The BRFSS is a state-based surveillance system administered by the Centers for
Disease Control and Prevention (CDC), providing random-digit-dialed telephone
interviews with non-institutionalized U.S.\ adults. All data are de-identified
and publicly available, with no ethics review required for secondary analysis of
publicly released administrative survey data.

\subsection{Cohort Selection}

We restricted the analytic cohort to respondents satisfying four criteria:
(i) interviewed in states administering the CIMEMLO1 optional module in 2023 or
2024; (ii) aged $\geq$50 years; (iii) providing a valid response to the SCD
outcome item; and (iv) not missing $>$80\% of predictor features. Figure 1
(cohort flow diagram) summarizes the selection process. From a raw pooled sample
of 3,968,799 BRFSS records, 2,544,716 were aged $\geq$50. After restricting to
CIMEMLO1-administering states and applying data quality filters, the final
analytic cohort comprised $\mathbf{N = 192{,}567}$ respondents (2023: $n = 121{,}970$;
2024: $n = 70{,}597$).

\subsection{Outcome Variable}

The SCD outcome was derived from the BRFSS item: \textit{``During the past
12 months, have you experienced confusion or memory loss that is happening
more often or getting worse?''} Affirmative responses were coded SCD-positive
($n = 32{,}655$; 17.0\%); negative responses SCD-negative ($n = 159{,}912$;
83.0\%). The SCD item has been part of the BRFSS Cognitive Decline module since
2011 with consistent wording, enabling valid pooling across survey years
\cite{Jessen2014}.

\subsection{Feature Engineering}

We constructed 33 candidate predictors from BRFSS items across five domains:
(i) \textit{Social isolation}: loneliness frequency (LONELYFREQ), social
support adequacy (SOCLSUPRT), frequency of social activities (SOCFREQ), living
arrangement (HLTHPLN1\_proxy), and emotional support availability (EMTSUPRT);
(ii) \textit{Mental health}: depression diagnosis (DEPRESSION), number of
poor mental health days in the past month (MENTHLTH), and anxiety proxy
(DECIDE); (iii) \textit{Physical health}: self-rated general health (GENHLTH),
number of poor physical health days (PHYSHLTH), arthritis diagnosis (ARTHRITIS),
smoking history (SMOKE100), and healthcare access (PERSDOC3);
(iv) \textit{Sociodemographic}: age group (AGE), sex (SEX), race/ethnicity
(RACE), education (EDUCA), income level (INCOME3), employment (EMPLOY1),
and marital status (MARITAL); and (v) \textit{Composite isolation score}: a
Multi-dimensional Social Isolation Index (MSII) derived from a weighted sum
across all available BRFSS social module items, with higher scores indicating
greater isolation (range 0--20; high isolation threshold $\geq$12.5).

Missing values were handled using scikit-learn's \texttt{SimpleImputer} with
median strategy for continuous features and most-frequent strategy for
categorical features. The survey year 2023 lacked SLEPTIM responses; after
imputation the feature matrix contained 31 active features (two sleep-related
items dropped at imputation). All features were encoded as numeric arrays
without standardization (tree models), with logistic regression using
min-max scaling applied separately to the training fold.

\subsection{Machine Learning Pipeline}

Five classifiers were trained and compared: logistic regression (LR),
random forest (RF), XGBoost~\cite{ChenGuestrin2016}, LightGBM~\cite{Ke2017},
and a stacking ensemble (LR meta-learner on LR, RF, XGBoost, LightGBM base
learners). Class imbalance (SCD+ : SCD- = 1 : 4.9) was addressed using
SMOTE oversampling applied exclusively within training folds to prevent
information leakage. Hyperparameters were tuned via 5-fold stratified
cross-validation on the 2023 development set using Bayesian optimization
(Optuna). The primary performance metric for hyperparameter selection was AUROC.
All pipelines were implemented in scikit-learn 1.5.2 \cite{Pedregosa2011},
xgboost 3.2.0, and lightgbm 4.7.0.

\subsection{Temporal Validation Strategy}

To simulate prospective deployment, we employed a \textit{temporal holdout}
validation: models were trained exclusively on 2023 BRFSS data ($n = 121{,}970$)
and evaluated on 2024 data ($n = 70{,}597$). This is an \textbf{internal
temporal holdout}, not an external validation; the two waves originate from the
same BRFSS survey program. Internal temporal validation guards against
contemporaneous data leakage and provides a realistic estimate of prospective
discrimination.

\subsection{Performance Evaluation}

Primary discrimination was assessed using the Area Under the Receiver Operating
Characteristic Curve (AUROC). Secondary metrics included the Area Under the
Precision-Recall Curve (AUPRC), Brier score, F1 score, sensitivity,
specificity, positive predictive value (PPV), negative predictive value (NPV),
and number needed to screen (NNS = $1/\text{PPV}$). The optimal decision
threshold was selected by maximizing the Youden index ($J = \text{sensitivity}
+ \text{specificity} - 1$) on the development set. DeLong's nonparametric method
\cite{DeLong1988} was used to compute 95\% confidence intervals for AUROC.
Model calibration was assessed visually via calibration plots and quantified by
the Brier score \cite{Steyerberg2010}.
To assess potential response-selection bias, a survey-weight
sensitivity analysis was conducted: AUROC was recomputed using BRFSS
complex survey weights (\texttt{LLCPWT}), following CDC analytic guidelines
for population-representative inference.
Clinical utility was assessed by decision curve analysis (DCA) \cite{Vickers2006},
computing net benefit for the model and treat-all/treat-none
benchmarks across the clinically plausible threshold range 0.05--0.50.

\subsection{SHAP Interpretability Analysis}

SHAP (SHapley Additive exPlanations) values \cite{LundbergLee2017} were
computed for XGBoost using the native \texttt{pred\_contribs=True} argument on
the XGBoost DMatrix interface (compatible with XGBoost 3.x / SHAP 0.49.1). A
random sample of $n = 5{,}000$ holdout observations was used for SHAP
computation to manage computational cost. Feature importance was ranked by mean
absolute SHAP value across the sample. SHAP beeswarm and dependence plots were
generated to characterize the direction and non-linearity of top-predictor effects.

\subsection{MSII Feature Contribution (Ablation Analysis)}

To quantify the incremental predictive contribution of the MSII features,
an ablation analysis was performed by retraining the full XGBoost pipeline
with all five MSII-derived features removed (\texttt{MSII\_COMPOSITE},
\texttt{MSII\_HIGH}, \texttt{MSII\_D1}--\texttt{MSII\_D3}).
AUROC and AUPRC were compared on the holdout set between the full model
and the MSII-ablated model. The Net Reclassification Index (NRI) was
computed to quantify reclassification improvement for events (SCD-positive)
and non-events (SCD-negative) separately \cite{Pencina2008}.

\subsection{Subgroup Analysis}

Pre-specified subgroup AUROCs were computed for nine strata: three age bands
(50--64, 65--79, $\geq$80), sex (male, female), MSII isolation level (high:
MSII $\geq$12.5 vs.\ low), depression status (diagnosed vs.\ not), and survey
year (2023--2024 pooled). No adjustment for multiple comparisons was made;
subgroup analyses are descriptive.

%% ── 3. Results ───────────────────────────────────────────────────────────────
\section{Results}

\subsection{Cohort Characteristics}

Table~\ref{tab:table1} presents participant characteristics by SCD status. The
analytic cohort comprised $192{,}567$ adults (mean age $66.8 \pm 9.2$ years;
55.2\% female). SCD was present in $32{,}655$ respondents (17.0\%; class ratio
4.9 : 1). Compared with SCD-negative adults, SCD-positive adults were slightly
older ($67.4 \pm 9.4$ vs.\ $66.7 \pm 9.1$ years), had higher mean MSII scores
($10.7 \pm 3.5$ vs.\ $8.8 \pm 3.0$), and were more than twice as likely to have
a depression diagnosis (40.0\% vs.\ 14.6\%). High social isolation (MSII
$\geq$12.5) was present in 36.4\% of SCD-positive vs.\ 23.7\% of SCD-negative
respondents. Education $\geq$some college was near-universal (94.5% overall)
with modest differences by SCD status (92.4\% SCD+ vs.\ 94.9\% SCD-).

\begin{table}[H]
\centering
\caption{Participant Characteristics by Subjective Cognitive Decline Status
(BRFSS 2023--2024, $N = 192{,}567$)}
\label{tab:table1}
\small
\begin{tabular}{L{5.5cm} C{2.8cm} C{2.8cm} C{2.8cm}}
\toprule
\textbf{Characteristic} &
\textbf{All participants} &
\textbf{SCD-positive} &
\textbf{SCD-negative} \\
 & ($N = 192{,}567$) & ($n = 32{,}655$; 17.0\%) & ($n = 159{,}912$; 83.0\%) \\
\midrule
\multicolumn{4}{l}{\textit{Demographics}} \\
Age, mean $\pm$ SD, years & $66.8 \pm 9.2$ & $67.4 \pm 9.4$ & $66.7 \pm 9.1$ \\
Female sex, $n$ (\%) & $106{,}368$ (55.2\%) & $18{,}401$ (56.3\%) & $87{,}967$ (55.0\%) \\
Male sex, $n$ (\%) & $86{,}199$ (44.8\%) & $14{,}254$ (43.7\%) & $71{,}945$ (45.0\%) \\
Education $\geq$ some college, $n$ (\%) & $181{,}226$ (94.5\%) & $30{,}106$ (92.4\%) & $151{,}120$ (94.9\%) \\
\midrule
\multicolumn{4}{l}{\textit{Social isolation}} \\
MSII score, mean $\pm$ SD & $9.1 \pm 3.1$ & $10.7 \pm 3.5$ & $8.8 \pm 3.0$ \\
High isolation (MSII $\geq 12.5$), $n$ (\%) & $49{,}730$ (25.8\%) & $11{,}881$ (36.4\%) & $37{,}849$ (23.7\%) \\
\midrule
\multicolumn{4}{l}{\textit{Mental health}} \\
Depression diagnosis, $n$ (\%) & $36{,}264$ (18.9\%) & $12{,}944$ (40.0\%) & $23{,}320$ (14.6\%) \\
\midrule
\multicolumn{4}{l}{\textit{Study design}} \\
Development set (2023 wave) & $121{,}970$ & $20{,}661$ & $101{,}309$ \\
Holdout set (2024 wave) & $70{,}597$ & $11{,}994$ & $58{,}603$ \\
\bottomrule
\end{tabular}
\begin{flushleft}
\small MSII = Multi-dimensional Social Isolation Index (range 0--20; high
threshold $\geq$12.5). SCD = subjective cognitive decline (BRFSS CIMEMLO1 item).
SD = standard deviation.
\end{flushleft}
\end{table}

\subsection{Model Comparison}

Table~\ref{tab:model_comparison} presents cross-validation and holdout
performance for all five classifiers. XGBoost achieved the highest holdout
AUROC (0.7678, DeLong 95\% CI 0.7595--0.7761) and was selected as the primary
model. LightGBM performed comparably (holdout AUROC 0.7669). The stacking ensemble underperformed all individual base models (holdout AUROC 0.7412). A post-hoc diagnostic confirmed this was driven by near-redundant base predictions: mean pairwise Pearson correlation among all five base-model holdout probability vectors was $r = 0.927$. Under such high inter-model correlation, the linear meta-learner's out-of-fold calibration surface becomes near-degenerate, and the stacking pipeline overfits the limited prediction-difference signal. Consistent with this, a simple average ensemble (unweighted mean of five base-model probability vectors) achieved holdout AUROC 0.7679 --- marginally exceeding XGBoost alone (0.7678) --- confirming that base-model diversity, not combination method, is the binding constraint on ensemble performance. Logistic regression performed
competitively given its parsimony (holdout AUROC 0.7572), confirming meaningful
linear signal in the feature space.

\begin{table}[H]
\centering
\caption{Model Performance --- 5-Fold Cross-Validation and 2024 Holdout Test Set
($N_{\text{holdout}} = 70{,}597$)}
\label{tab:model_comparison}
\small
\begin{tabular}{L{3.5cm} C{2.4cm} C{2.4cm} C{2.4cm} C{2.4cm}}
\toprule
\textbf{Model} &
\textbf{CV AUROC} &
\textbf{CV AUROC 95\% CI} &
\textbf{Holdout AUROC} &
\textbf{Holdout AUROC 95\% CI} \\
\midrule
Logistic Regression & 0.7475 & --- & 0.7572 & --- \\
Random Forest & 0.7514 & --- & 0.7598 & --- \\
\textbf{XGBoost}$^{*}$ & \textbf{0.7586} & \textbf{[0.7553, 0.7608]} & \textbf{0.7678} & \textbf{[0.7595, 0.7761]} \\
LightGBM & 0.7576 & --- & 0.7669 & --- \\
Stacking Ensemble & 0.7314 & --- & 0.7412 & --- \\
\bottomrule
\end{tabular}
\begin{flushleft}
\small $^{*}$Primary model. CV AUROC is the mean of 5-fold stratified
cross-validation on the 2023 development set ($n = 121{,}970$). Holdout AUROC
95\% CI computed by DeLong's method \cite{DeLong1988}. CI not reported for
non-primary models; primary model CI shown to three decimal places.
\end{flushleft}
\end{table}

\subsection{Primary Model Performance: XGBoost}

On the 2024 holdout test set, XGBoost achieved:
\begin{itemize}[noitemsep]
  \item AUROC: 0.7678 (95\% CI 0.7595--0.7761)
  \item AUPRC: 0.4339
  \item Brier score: 0.1193
  \item F1 score: 0.4366
\end{itemize}

At the Youden-optimal decision threshold of 0.1697:
\begin{itemize}[noitemsep]
  \item Sensitivity: 0.72 (72\% of true SCD cases detected)
  \item Specificity: 0.68 (68\% of true SCD-negative cases correctly excluded)
  \item Positive predictive value (PPV): 0.31 (31\% of model-flagged individuals have SCD)
  \item Negative predictive value (NPV): 0.92 (92\% of those below threshold are SCD-negative)
  \item Number needed to screen (NNS): $1/0.3133 \approx 3.2$
\end{itemize}

\noindent\textbf{Survey-weight sensitivity analysis.} Survey-weighted AUROC (using BRFSS \texttt{LLCPWT} complex survey weights) was 0.7678, identical to the unweighted estimate ($\Delta$AUROC~=~0.0000). This confirms that BRFSS respondents in CIMEMLO1-administering states are not differentially over- or under-represented in a way that biases discrimination estimates, supporting generalizability to the U.S.\ older-adult population.

The high NPV (0.92) is particularly clinically relevant: individuals below the
model threshold have a 92\% probability of being SCD-negative, supporting the
model's utility as a first-tier rule-out instrument. The NNS of 3.2 indicates
that approximately 3 individuals need to be flagged by the model to identify
one true SCD case --- favourably comparable to breast cancer screening NNS
benchmarks.

\subsection{Model Calibration}

Calibration curves (Figure~3) showed that XGBoost
produced well-calibrated probability estimates in the 0.10--0.30 predicted
probability range, with minor over-prediction above 0.35 (consistent with
high-risk subgroups having elevated true prevalence). The Brier score of 0.1193
compares favorably to the naive no-information Brier score of
$0.17 \times 0.83 = 0.141$, confirming model skill beyond prevalence-only
prediction.

\subsection{Decision Curve Analysis}

Decision curve analysis (Figure~4) demonstrated that the XGBoost model
provided net clinical benefit across the 0.10--0.35 predicted-probability
range. At the Youden-optimal threshold of 0.1697, the model net benefit was
0.0673, while the treat-all strategy yielded a net benefit of $-$0.0001 ---
effectively zero --- confirming that the model substantially outperforms
universal screening for the relevant moderate-risk stratum. The treat-none
strategy had net benefit 0 by definition. These results support clinical
deployment of the model as a threshold-guided screening instrument.

\subsection{SHAP Feature Importance}

Figure~5 presents the SHAP beeswarm plot for XGBoost on
the $n = 5{,}000$ holdout sample. The top 10 predictors by mean $|$SHAP$|$ were:

\begin{enumerate}[noitemsep]
  \item \textbf{DEPRESSION} (depression diagnosis) --- highest mean $|$SHAP$|$;
        strong positive association (diagnosed = higher SCD risk)
  \item \textbf{MENTHLTH} (poor mental health days, 0--30) --- continuous
        dose-response (more poor days = higher SCD risk)
  \item \textbf{GENHLTH} (self-rated general health, 1--5) --- poor health
        strongly positive
  \item \textbf{ARTHRITIS} (arthritis diagnosis) --- moderate positive
        association
  \item \textbf{AGE} (age group) --- older age modestly associated with higher
        SCD risk
  \item \textbf{SMOKE100} (ever smoked $\geq$100 cigarettes) --- positive
        association
  \item \textbf{MSII\_D2} (MSII isolation domain 2 component) --- positive;
        7th-ranked social isolation signal
  \item \textbf{POORHLTH} (days unable to perform activities due to poor health)
        --- positive
  \item \textbf{PERSDOC3} (has a personal doctor) --- inverse association
        (having a doctor = lower SCD risk; possibly reflecting healthcare
        engagement)
  \item \textbf{SEX} (biological sex) --- modest; female sex shows slightly
        higher SCD probability in SHAP space
\end{enumerate}

The MSII component (MSII\_D2) ranked 7th, confirming that multi-dimensional
social isolation carries incremental predictive signal beyond depression and
general health status alone. SHAP dependence plots (Supplementary Figure~S1)
revealed a non-linear interaction between DEPRESSION and MENTHLTH: the SHAP
contribution of depression diagnosis was amplified among respondents with
$\geq$15 poor mental health days per month.

\subsection{MSII Feature Contribution (Ablation Analysis)}

Removing all five MSII features reduced holdout AUROC from 0.7678 to 0.7615
($\Delta$AUROC~=~$+$0.0063) and AUPRC from 0.4339 to 0.4264
($\Delta$AUPRC~=~$+$0.0075; Table~\ref{tab:ablation}). The NRI confirmed
meaningful reclassification improvement: NRI\textsubscript{events}~=~0.034
(3.4\% of SCD-positive individuals moved to a higher-risk decile),
NRI\textsubscript{non-events}~=~0.184, NRI\textsubscript{total}~=~0.218.
Although the absolute AUROC gain is modest, the NRI total of 0.218 indicates
that the MSII features provide non-trivial incremental reclassification benefit
for both cases and non-cases, supporting their retention in the final model.

\begin{table}[H]
\centering
\caption{MSII Feature Contribution: Ablation Analysis on 2024 Holdout Set}
\label{tab:ablation}
\small
\begin{tabular}{L{5.5cm} C{2.5cm} C{2.5cm} C{2cm}}
\toprule
\textbf{Metric} & \textbf{Full Model} & \textbf{Without MSII} & \textbf{$\Delta$} \\
\midrule
AUROC (95\% CI) & 0.7678 & 0.7615 & $+$0.0063 \\
AUPRC & 0.4339 & 0.4264 & $+$0.0075 \\
NRI\textsubscript{events} & --- & --- & $+$0.034 \\
NRI\textsubscript{non-events} & --- & --- & $+$0.184 \\
NRI\textsubscript{total} & --- & --- & $+$0.218 \\
\bottomrule
\end{tabular}
\begin{flushleft}
\small MSII = Multi-dimensional Social Isolation Index. Full model includes
MSII\_COMPOSITE, MSII\_HIGH, and three domain sub-scores. NRI = Net
Reclassification Index \cite{Pencina2008}; positive values indicate improvement
with MSII features included. Holdout $N = 70{,}597$.
\end{flushleft}
\end{table}

\subsection{Subgroup Analysis}

Table~\ref{tab:subgroup} presents pre-specified subgroup AUROCs with
bootstrap 95\% confidence intervals (500 resamples) on the 2024 holdout set.
Model discrimination was strongest in the 50--64 age group (AUROC 0.8166,
95\% CI 0.8099--0.8230), and declined monotonically with advancing age
(65--79: 0.7482 [0.7404--0.7557]; $\geq$80: 0.6734 [0.6593--0.6863]).
Discrimination was similar between males (0.7698 [0.7638--0.7767]) and
females (0.7665 [0.7603--0.7734]). MSII-high and MSII-low subgroups performed
similarly (0.7562 vs.\ 0.7607). The depression subgroup showed lower AUROC
(0.7148 [0.7061--0.7235]), consistent with depression-mediated SCD risk
reducing the residual discriminatory signal from social isolation features.
High between-subgroup heterogeneity was observed (Cochran's $Q = 202.6$,
$I^{2} = 96.5\%$, $p < 0.001$), driven primarily by the age-gradient
effect, suggesting that age-stratified deployment thresholds warrant
prospective investigation.

\begin{table}[H]
\centering
\caption{Pre-Specified Subgroup Analysis on the 2024 Holdout Test Set (Bootstrap 95\% CIs, $B = 500$ resamples)}
\label{tab:subgroup}
\small
\begin{tabular}{L{5cm} C{2cm} C{2.5cm} C{3.5cm}}
\toprule
\textbf{Subgroup} & \textbf{$n$} & \textbf{AUROC} & \textbf{95\% CI} \\
\midrule
\multicolumn{4}{l}{\textit{Age group}} \\
\quad Age 50--64 years & 27,877 & \textbf{0.8166} & [0.8099, 0.8230] \\
\quad Age 65--79 years & 30,603 & 0.7482 & [0.7404, 0.7557] \\
\quad Age $\geq$80 years & 10,091 & 0.6734 & [0.6593, 0.6863] \\
\midrule
\multicolumn{4}{l}{\textit{Sex}} \\
\quad Male & 31,877 & 0.7698 & [0.7638, 0.7767] \\
\quad Female & 38,720 & 0.7665 & [0.7603, 0.7734] \\
\midrule
\multicolumn{4}{l}{\textit{Social isolation level}} \\
\quad MSII high ($\geq$12.5) & 16,382 & 0.7562 & [0.7477, 0.7646] \\
\quad MSII low ($<$12.5) & 54,215 & 0.7607 & [0.7549, 0.7664] \\
\midrule
\multicolumn{4}{l}{\textit{Mental health comorbidity}} \\
\quad Depression diagnosis & 14,010 & 0.7148 & [0.7061, 0.7235] \\
\midrule
\multicolumn{4}{l}{\textit{Heterogeneity (all subgroups)}} \\
\quad Cochran's $Q$ & \multicolumn{3}{l}{202.6, $df = 7$, $p < 0.001$} \\
\quad $I^{2}$ & \multicolumn{3}{l}{96.5\% (very high)} \\
\bottomrule
\end{tabular}
\begin{flushleft}
\small All $n$ values are holdout-only (SURVEY\_YEAR = 2024, $N_{\text{hold}} = 70{,}597$).
AUROC = Area Under the Receiver Operating Characteristic Curve;
CI = confidence interval from 500 stratified bootstrap resamples.
MSII = Multi-dimensional Social Isolation Index.
$I^{2}$ heterogeneity driven primarily by the age gradient.
\end{flushleft}
\end{table}

%% ── 4. Discussion ────────────────────────────────────────────────────────────
\section{Discussion}

\subsection{Principal Findings}

This study presents the largest U.S.\ ML analysis ($N = 192{,}567$) linking
multi-dimensional social isolation features to subjective cognitive decline
in community-dwelling adults aged $\geq$50. XGBoost achieved a holdout AUROC
of 0.7678 --- the highest among five models --- with a clinically useful NPV of
0.92 and NNS of 3.2. These results establish that a no-biomarker, no-clinical-visit
screening instrument, deployable from routine population health survey responses,
can detect SCD with moderate-to-good discrimination.

The pre-specified hypothesis from the study protocol was AUROC $\geq$ 0.80
(benchmarked against Fasokun et al. \cite{Fasokun2026} and Alam et al.\
\cite{Alam2026}). The achieved AUROC of 0.7678 fell below this target.
Transparent reporting of this gap is essential: the achieved performance
reflects the inherent difficulty of predicting SCD from cross-sectional survey
data in a large, heterogeneous national sample. The models in Ishikawa et al.
\cite{Ishikawa2025} and Alam et al.\ \cite{Alam2026} achieved AUROC $\sim$0.82
in substantially smaller, potentially less heterogeneous samples. Discrimination
typically decreases as sample size and population heterogeneity increase,
which partially explains the gap. Importantly, the NPV of 0.92 and NNS of 3.2
remain clinically meaningful for a first-tier population screening application.

\subsection{Comparison with Prior Literature}

The closest comparable work is Nguyen et al.\ \cite{Nguyen2026}, who applied
gradient-boosted trees to pooled BRFSS data ($N = 298{,}944$; AUROC
$\approx$0.76--0.77), establishing a BRFSS-specific performance ceiling
of approximately AUROC 0.77. The present study achieves consistent performance
(AUROC 0.7678) while contributing a distinct thematic contribution: the first
BRFSS ML analysis to operationalize the MSII as the primary exposure
framework, quantify its incremental contribution via ablation and NRI, and
report a fully actionable screening threshold with NNS. Among clinical-cohort
studies, Alam et al.\ \cite{Alam2026} ($N = 5{,}984$) and Ishikawa et al.\
\cite{Ishikawa2025} ($N = 4{,}586$) achieved AUROC $\sim$0.82 using richer
clinical features at sample sizes more than 30-fold smaller.
The only prior BRFSS-based association study (Fasokun et al.\ \cite{Fasokun2026},
$N = 86{,}520$) reported no ML model or AUROC.

The critical methodological distinction applies to all three BRFSS comparators:
Fasokun et al.\ addressed \textit{``does loneliness correlate with SCD?''};
Nguyen et al.\ addressed \textit{``can gradient boosting predict SCD from
BRFSS survey data?''}; the present study addresses \textit{``can ML detect SCD
using specifically multi-dimensional social isolation features, with actionable
screening metrics?''}. These are distinct scientific questions. Fasokun et al.\ addressed the question ``\textit{does
loneliness correlate with SCD?}''; the present study addresses ``\textit{can ML
predict who will screen positive for SCD using social isolation features, with
sufficient accuracy for population deployment?}'' These are fundamentally
different scientific questions with different clinical implications.

Chinese cohort studies \cite{Wang2022,Hu2021} and Korean data \cite{Oh2024}
achieved AUROCs of 0.81--0.83 but used different survey instruments, healthcare
systems, and social constructs, limiting direct comparability. The U.K.\ ELSA
study \cite{Ahmadzadeh2023} achieved AUROC 0.73 in $N = 2{,}733$, closer to
the current result and likely reflecting similar challenges in nationally
representative heterogeneous cohorts.

\subsection{Clinical and Public Health Interpretation}

The dominance of depression and mental health features (DEPRESSION, MENTHLTH)
in SHAP rankings confirms a well-established clinical co-occurrence: depression
is both a risk factor for and a symptom companion to cognitive decline. The
MSII isolation composite ranked 7th, providing incremental signal but not
dominating prediction. This finding does not diminish the relevance of social
isolation; rather, it suggests that in BRFSS data, the \textit{experienced
consequence} of isolation (depression, poor mental health) is more immediately
predictive of SCD than the isolation exposure itself. Social isolation may
operate as an \textit{upstream determinant} whose pathway to SCD is mediated
through depression and general health deterioration.

The subgroup finding that the 50--64 age band achieves AUROC 0.82 is clinically
meaningful: this age group is likely in the pre-clinical SCD window when
preventive intervention --- social prescribing, depression treatment, community
connection programs --- can realistically alter cognitive trajectories. The
lower AUROC in adults $\geq$80 (0.70) reflects multiple concurrent etiologies
of cognitive difficulty at advanced age, reducing the discriminatory contribution
of social isolation features.

The NNS of 3.2 compares favorably with population screening benchmarks.
Specifically, for a condition with 17\% prevalence and the social isolation risk
pathway, a first-tier screen that identifies one true case for every 3.2 flagged
individuals is competitive with primary care depression screening (NNS typically
3--5 for PHQ-based tools at similar thresholds). Implementing this model as a
population-level BRFSS score could guide states toward targeting social
prescribing resources at high-risk adults before functional impairment emerges.

\subsection{Limitations}

Several limitations must be acknowledged. First, this is a \textbf{cross-sectional}
study: exposure and outcome are measured simultaneously, precluding causal
inference. The temporal direction of the social isolation -- SCD relationship
cannot be established from BRFSS data alone; reverse causation (SCD causing
social withdrawal) is plausible. Second, the SCD outcome is \textbf{self-reported
and single-item}. It captures perceived cognitive change but lacks clinical
validation against neuropsychological testing or biomarker confirmation.
Approximately 40--50\% of BRFSS SCD-positive respondents may not meet clinical
criteria for mild cognitive impairment. Third, the \textbf{temporal holdout is
internal}: 2023 training to 2024 testing involves the same BRFSS infrastructure,
questionnaire, and sample frame. True external validation in an independent
cohort (e.g., HRS, NHATS) is needed before clinical deployment.

Fourth, the achieved AUROC of 0.7678 \textbf{falls below the pre-specified
target of 0.80}. Feature availability in BRFSS is limited by the optional-module
structure: not all states administer all BRFSS social modules in both years.
Enrichment with clinical variables (biomarkers, neuropsychological test scores)
would likely improve discrimination but is not available in BRFSS. Fifth,
BRFSS telephone sampling \textbf{underrepresents} the most socially isolated
adults (those without telephone access or too cognitively impaired to complete
the survey), introducing ascertainment bias that may attenuate true AUROC.
Sixth, \textbf{race/ethnicity and geographic region} were included as features
but BRFSS respondent misclassification and varying state participation in
CIMEMLO1 may limit generalizability across demographic subgroups and states.

\subsection{Future Directions}

Several extensions are warranted. Longitudinal BRFSS data linkage or integration
with Medicare claims records could convert this cross-sectional design to a
prospective framework for incident SCD prediction. Federated learning approaches
across state BRFSS datasets could improve model generalizability while preserving
data privacy \cite{Alam2026}. External validation in NHATS or HRS is a priority.
Finally, a clinical decision support tool embedding the XGBoost score in an
electronic health record prompt could enable prospective evaluation of NNS in
real primary care settings.

%% ── 5. Conclusions ───────────────────────────────────────────────────────────
\section{Conclusions}

An XGBoost classifier trained on routinely collected BRFSS social and health
survey data discriminates subjective cognitive decline with AUROC 0.7678
(95\% CI 0.7595--0.7761) in a holdout cohort of 70,597 older U.S.\ adults,
delivering an NPV of 0.92 and NNS of 3.2 at the Youden-optimal threshold. To our knowledge, this is among the largest BRFSS ML studies focused on
multi-dimensional social isolation as the primary exposure framework for
SCD screening ($N = 192{,}567$), with AUROC consistent with the concurrent
BRFSS benchmark \cite{Nguyen2026}. Decision curve analysis confirms net
clinical benefit at the deployment threshold; MSII features provide
incremental reclassification benefit (NRI\textsubscript{total}~=~0.218)
supporting their inclusion.
SHAP analysis reveals that depression and mental health days dominate prediction,
with the MSII social isolation composite contributing incremental 7th-rank signal,
suggesting that the social pathway to SCD operates partly through mental health
as an intermediate. The model falls below the pre-specified AUROC target of 0.80,
consistent with the challenges of nationally representative heterogeneous survey
data, but maintains clinically useful NPV and NNS. Before deployment, external
validation and a prospective clinical decision study are required.

%% ── Declarations ─────────────────────────────────────────────────────────────
\section*{Declarations}

\subsection*{Ethics Approval and Consent to Participate}
This study used publicly available, de-identified BRFSS data released by the
CDC. No human subjects research review was required for secondary analysis of
publicly released, non-identifiable administrative data.

\subsection*{Availability of Data and Materials}
BRFSS annual data files are publicly available at
\url{https://www.cdc.gov/brfss/annual\_data/annual\_data.htm}. All analysis
scripts are available at [repository URL pending publication].

\subsection*{Competing Interests}
The author declares no competing interests.

\subsection*{Funding}
No external funding was received for this study.

\subsection*{Authors' Contributions}
[Author]: Conceptualization, methodology, software, formal analysis, data
curation, writing (original draft, review and editing).

\subsection*{Acknowledgments}
The author acknowledges the CDC and all BRFSS state health departments for
collecting and publicly releasing these data. The author used SciSpace AI
Research Agent (SciSpace, 2026) as a generative AI assistant for analytical
pipeline construction and manuscript drafting support; see the AI Usage
Declaration on the title page.

%% ── Abbreviations ────────────────────────────────────────────────────────────
\section*{Abbreviations}
\begin{tabular}{ll}
ADRD & Alzheimer's disease and related dementias \\
AUPRC & Area Under the Precision-Recall Curve \\
AUROC & Area Under the Receiver Operating Characteristic Curve \\
BRFSS & Behavioral Risk Factor Surveillance System \\
CDC & Centers for Disease Control and Prevention \\
CIMEMLO1 & Cognitive Impairment and Memory Loss optional module 1 \\
GBM & Gradient Boosted Machine \\
LGB & LightGBM \\
LR & Logistic Regression \\
ML & Machine Learning \\
MCI & Mild Cognitive Impairment \\
MSII & Multi-dimensional Social Isolation Index \\
NNS & Number Needed to Screen \\
NPV & Negative Predictive Value \\
PPV & Positive Predictive Value \\
RF & Random Forest \\
SCD & Subjective Cognitive Decline \\
SDoH & Social Determinants of Health \\
SHAP & SHapley Additive exPlanations \\
XGB & XGBoost \\
\end{tabular}

%% ── References ───────────────────────────────────────────────────────────────
\bibliographystyle{unsrtnat}
\bibliography{references}

%% ── Figure Legends ───────────────────────────────────────────────────────────
\newpage
\section*{Figure Legends}

\noindent\textbf{Figure 1. Cohort Flow Diagram.}
Selection of the analytic cohort from BRFSS 2023--2024. Starting from
3,968,799 total BRFSS records, sequential filters for age $\geq$50,
CIMEMLO1 module completion, and data quality yielded $N = 192{,}567$
for analysis ($n = 121{,}970$ in the 2023 development set;
$n = 70{,}597$ in the 2024 holdout test set).

\vspace{0.5em}

\noindent\textbf{Figure 2. ROC Curves for All Five Classifiers on the
2024 Holdout Test Set.}
Receiver Operating Characteristic (ROC) curves for logistic regression (AUROC
0.7572), random forest (0.7598), XGBoost (0.7678, primary model), LightGBM
(0.7669), and stacking ensemble (0.7412). The diagonal dashed line represents
chance performance. DeLong 95\% CI for XGBoost: [0.7595, 0.7761].

\vspace{0.5em}

\noindent\textbf{Figure 3. Calibration Curve --- XGBoost on the 2024 Holdout Test Set.}
Observed event rate (binned, 20 bins) versus mean predicted probability for
XGBoost on the $N = 70{,}597$ holdout set. The diagonal dashed line represents
perfect calibration. Calibration slope~=~1.030; intercept~=~$-$0.072;
Brier score~=~0.1193. Minor over-prediction is observed above predicted
probability 0.35, consistent with elevated true prevalence in high-risk strata.

\vspace{0.5em}

\noindent\textbf{Figure 4. Decision Curve Analysis --- XGBoost vs.\ Benchmarks.}
Net benefit curves for the XGBoost model, the treat-all strategy,
and the treat-none strategy across the threshold probability range
0.05--0.50. At the Youden-optimal threshold of 0.1697, XGBoost net
benefit~=~0.0673; treat-all net benefit~=~$-$0.0001. The model provides
positive net benefit over both benchmarks across the 0.10--0.35 range,
confirming clinical utility of the screening threshold \cite{Vickers2006}.

\vspace{0.5em}

\noindent\textbf{Figure 5. SHAP Beeswarm Plot --- XGBoost Top 20 Features
($n = 5{,}000$ Holdout Sample).}
Each point represents one respondent; x-axis shows SHAP value (impact on
log-odds of SCD-positive prediction); color encodes feature value (red =
high, blue = low). Features ranked by mean $|$SHAP$|$ descending. Top feature
DEPRESSION shows the largest spread; MSII\_D2 (social isolation composite)
ranks 7th.

\section*{Supplementary Material}

\noindent\textbf{Table S1. Full Hyperparameter Configurations.}
[Provided as supplementary file pending repository upload.]

\noindent\textbf{Table S2. SHAP Values for All 31 Features.}
[Provided as supplementary file pending repository upload.]

\noindent\textbf{Table S3. Calibration Metrics by Model.}
[Provided as supplementary file pending repository upload.]

\noindent\textbf{Table S4. TRIPOD+AI Checklist.}
Completed TRIPOD+AI reporting checklist \cite{Collins2024} for this
prediction model study. [Provided as supplementary file pending
repository upload.]

\end{document}
